\clearpage
\begin{examplebox}
\textbf{\textcolor{CExample}{例1：直接求两条距离——右端点检查符号}}

\textbf{题目：} 椭圆 $x^2/25+y^2/16=1$ 上点 $P$ 的横坐标为 $5$，求两条焦半径。

\textbf{解：} $a=5,b=4,c=3,e=3/5$，故
\[
r_1=5+\frac35\cdot5=\boxed8,\qquad r_2=5-\frac35\cdot5=\boxed2.
\]
\textbf{回验：} 此时 $P=(5,0)$，到 $(-3,0)$、$(3,0)$ 的距离确为 $8,2$，和为 $10$。
这正是理解页三态图的右端状态。

\medskip
\textbf{\textcolor{CExample}{例2：由右焦半径反求横坐标}}

\textbf{题目：} 同一椭圆上 $|PF_2|=4$，求横坐标。

\textbf{解：} 右焦半径用减号，
\[
5-\frac35x_0=4\quad\Longrightarrow\quad \boxed{x_0=\frac53}.
\]
\textbf{回验：} 代入椭圆得 $y_0=\pm\frac{8\sqrt2}{3}$，确有上下两个点；
到右焦点的距离平方为 $(5/3-3)^2+128/9=16$，故距离为 $4$。

\medskip
\textbf{\textcolor{CExample}{例3：两条焦半径的乘积最大值}}

\textbf{题目：} 同一椭圆上求 $r_1r_2$ 的最大值。

\textbf{解：} $x_0\in[-5,5]$，所以
\[
r_1r_2=\left(5+\frac35x_0\right)\left(5-\frac35x_0\right)
=25-\frac9{25}x_0^2\le\boxed{25}.
\]
等号当且仅当 $x_0=0$，即 $P=(0,\pm4)$。
\textbf{回验：} 此时两条距离均为 $5$，乘积为 $25$；对应三态图的中间状态。
\end{examplebox}
\clearpage
\begin{examplebox}
\textbf{\textcolor{CExample}{例4：由两焦半径之差反求横坐标}}

\textbf{题目：} 椭圆 $x^2/25+y^2/16=1$ 上一点满足 $r_1-r_2=12/5$，求横坐标。

\textbf{识别：} 第一定义给和 $2a$，平方差法还给出 $r_1-r_2=2ex_0$。
代 $e=3/5$，
\[
\frac{12}{5}=2\cdot\frac35x_0\quad\Longrightarrow\quad\boxed{x_0=2}.
\]
\textbf{回验：} $P=(2,\pm4\sqrt{21}/5)$ 在椭圆上；
$r_1=31/5,r_2=19/5$，差为 $12/5$、和为 $10$。
这也是勾股构造图的左侧状态。

\medskip
\textbf{\textcolor{CExample}{例5：焦点在纵轴——下加上减}}

\textbf{题目：} 椭圆 $x^2/9+y^2/25=1$ 的下、上焦点分别为 $F_1,F_2$。
点 $P$ 的纵坐标 $y_0=5/2$，求两条焦半径。

\textbf{解：} $a=5,b=3,c=4,e=4/5$。焦点在纵轴，用 $y_0$：
\[
r_1=5+\frac45\cdot\frac52=\boxed7,\qquad
r_2=5-\frac45\cdot\frac52=\boxed3.
\]
\textbf{回验：} $P=(\pm3\sqrt3/2,5/2)$ 在椭圆上；
到下焦点 $(0,-4)$ 的距离平方为 $27/4+169/4=49$，
到上焦点 $(0,4)$ 的距离平方为 $27/4+9/4=9$，因此为 $7,3$，且和为 $10$。
\end{examplebox}
